\begin{proof}[Proof of \cref{obj:lem:class-compact-causal-range}]
Fix \(P\in\mathcal C\).  By the defining property of \(\mathcal C\), \(P\) satisfies all clauses of the transported model class \(\mathcal P_n\) at the index \(n\).  In particular, it is an admissible input for the source-observation facts of \cref{obj:lem:source-observation-facts}. % lean: ScoreRiskClassAtoms.toTransportedIVClass

First supply the side conditions required by \cref{obj:prop:compact-causal-range}.  The source-observation facts give, for every measurable \(A\subseteq\mathcal X\),
\[
\int_{\{X\in A\}}H(O^S)Y\,dP_S
=
\int_A\Delta_Y(x)\,dP_S^X(x),
\qquad
\int_{\{X\in A\}}H(O^S)D\,dP_S
=
\int_A\Delta_D(x)\,dP_S^X(x),
\]
and also
\[
\Delta_Y(X)\in[-1,1],
\qquad
\Delta_D(X)\in[0,1],
\qquad P_S^X\text{-a.s.}
\]
Consequently
\[
\Delta_Y\in L^1(P_S^X),
\qquad
\Delta_D\in L^1(P_S^X).
\]
% lean: transportedIVClass_contrast_integrable

Now apply \cref{obj:prop:compact-causal-range} to this \(P\) with those two integrability facts.  It gives the observed source contrast representations
\[
\int_{\{X\in A\}}H(O^S)Y\,dP_S
=
\int_A\Delta_Y(x)\,dP_S^X(x),
\qquad
\int_{\{X\in A\}}H(O^S)D\,dP_S
=
\int_A\Delta_D(x)\,dP_S^X(x)
\]
for every measurable \(A\), and the transported target identities
\[
\int w(x)\Delta_Y(x)\,dP_S^X(x)=\mu_{Y,n},
\qquad
\int w(x)\Delta_D(x)\,dP_S^X(x)=\mu_n.
\]
% lean: compact_causal_range

The same application identifies the Wald ratio with the target complier contrast:
\[
\frac{\mu_{Y,n}}{\mu_n}
=
\theta_T
=
\mathbb E_{P_n^F}\!\left[
Y(1)-Y(0)\mid D(1)=1,\ D(0)=0,\ S=0
\right],
\]
where
\[
\mu_{Y,n}
=
\int\bigl(Y(1)-Y(0)\bigr)\mathbf 1\{D(1)=1,\ D(0)=0\}\,
dP_n^F(\,\cdot\mid S=0),
\qquad
\mu_n=P_n^F\{D(1)=1,\ D(0)=0\mid S=0\}.
\]
Since \(Y(0),Y(1)\in[0,1]\) on the full-data support and the target complier event has positive probability by target-complier positivity, this conditional contrast lies in
\[
\Theta=[-1,1].
\]
% lean: compact_causal_range

Thus all conclusions of \cref{obj:prop:compact-causal-range} hold for the fixed \(P\in\mathcal C\).  Since \(P\) was arbitrary, they hold for every member of \(\mathcal C\). % lean: scoreRiskClass_compact_causal_range
\end{proof}
